\section{Content-Typed Endpoints}
\label{sec:content}

Three endpoints look inside a file: inline image preview, inline PDF, and
archive listing.  Each is defined by a predicate on the file's \emph{bytes},
never its name, and each factors through $\Open$, so Laws~\ref{law:absent}
and~\ref{law:indist} apply to them unchanged.  A fourth, the Quick Look picture,
hands bytes to another program; Law~\ref{law:quicklook} states what it may be
given.

\begin{definition}[Sniff]
$\Sniff_{\mathrm{img}}$ accepts a byte sequence that begins with the PNG, JPEG,
GIF, or WebP signature; $\Sniff_{\mathrm{pdf}}$ accepts one that begins with
\code{\%PDF-}; $\Sniff_{\mathrm{arc}}$ accepts the start of a zip, a tar
(\code{ustar} at offset 257), or a gzip stream that contains a tar.
\end{definition}

\begin{law}[Content decides]
\label{law:sniff}
The endpoint succeeds only if the corresponding sniff accepts the file's leading
bytes; the file name is never consulted.  A file named \code{a.png},
\code{a.pdf} or \code{a.zip} with other contents is refused with 415, and a file
with a different name but the right contents is accepted.  SVG and HTML are never
accepted.
\Tests{guard/m2_test.go (OpenImage); guard/pdf_test.go; guard/archive_test.go; server/pdf_test.go}
\end{law}

\begin{law}[Cloud-only files are never read]
\label{law:dataless}
For a file whose contents are stored only in the cloud, no endpoint reads it
implicitly: the guard reports it as not downloaded before any read.
\Tests{guard/dataless tests; e2e (hover peek)}
\end{law}

\begin{law}[Archive listing discloses names only]
\label{law:archive-names}
The archive listing contains, for each member, its name, size, type and date,
and never any member data; names are displayed as text with control characters,
Unicode formatting characters and invalid UTF-8 replaced, are cut to a bounded
prefix (4{,}096 bytes, then an ellipsis), and are never used as paths.
\Tests{guard/archive_test.go::TestListArchiveNamesAreDefanged; server/archive_test.go}
\end{law}

\begin{law}[Archive listing is bounded]
\label{law:archive-bounded}
For any file, listing examines at most $E$ tar entries or zip directory records,
streams at most $D$ decompressed bytes of a compressed tar, reads at most $B$
bytes of a zip directory, returns at most $S$ entries, and reports (as
\emph{incomplete}) when a scan limit stopped it.  A \emph{valid} archive that
exceeds a scan limit is listed up to the limit and marked incomplete; it is not
refused.  A zip whose directory was read in full is refused if its record count
disagrees with the archive's end record, which can be checked only then.  The
byte bound $B$ is checked between records, so the last record may cross it by at
most one record's length, and the work stays bounded.  A plain tar is read by
skipping each member's body, so its cost does not grow with member sizes.
\Tests{guard/archive\_test.go::TestListArchiveIsBounded, TestListArchiveStopsAtTheDecompressionLimit, TestZipEntryLimitHoldsForValidArchives, TestZipDirectoryByteBudget, TestZipWithAMisleadingCountIsRefused, TestTarGzLimitAtAHeaderBoundaryIsIncomplete, TestPlainTarListingStillSeeksPastMemberBodies}
\end{law}

\begin{law}[Inline PDF is framed only by the app]
\label{law:pdf-frame}
The PDF response carries \code{X-Frame-Options: SAMEORIGIN} and
\code{frame-ancestors 'self'}, the app shell carries \code{DENY} and
\code{frame-ancestors 'none'}, and the Markdown frame document carries its own
policy that admits exactly its one script and one style by hash.
\Tests{server/pdf_test.go; server/mdframe_test.go; web/mdframe_test.go; e2e PDF and Markdown tests}
\end{law}

\begin{law}[Quick Look sees only what $\Open$ serves]
\label{law:quicklook}
Let $C(p)$ be what Quick Look is given for a request on $p$.  $C(p)$ is defined
only if $\Open(p)$ succeeds, the real path has an iWork or Office extension, and
the file is not cloud-only.  For a file, $C(p)$ is a new file in a private folder
whose bytes are those read from the descriptor $\Open(p)$ returned.  For a
package, $C(p)$ is the set of pairs $(r, b)$ such that $\Open(p/r)$ succeeds on a
regular file whose real path lies under the real path of $p$, and $b$ is the bytes
read from that descriptor.  In both cases $C(p)$ carries no extended attributes and
no resource fork, and it is never the path $p$ itself.  Hence every byte Quick Look
reads is a byte some request could already read through $\Open$, and
Law~\ref{law:absent} extends to it: no denied file and nothing outside the roots
reaches it, whatever $p$ names (an alias, a link, a package with hidden parts).
\Tests{server/quicklook_test.go (private copy, no extended attributes, package parts, a real Finder alias); server/fuzz_test.go::FuzzPathParameter (stand-in Quick Look echoes its input)}
\end{law}

\begin{boundary}[What Quick Look does with the copy]
The law bounds Quick Look's \emph{input}, not its behavior: a flaw in Apple's
generators, or a document that refers to something outside itself, is not
modeled.  The independent review found no generator that loads local files named
inside a document.
\end{boundary}

\begin{boundary}[PDF sandboxing]
A CSP \code{sandbox} cannot be applied to a PDF, because browsers then refuse to
display it.  The protection is instead the byte check, the framing headers, and
the browser's own isolated PDF viewer.
\end{boundary}
